\documentclass[11pt]{article}
\usepackage[a4paper,margin=29mm]{geometry}
\usepackage{amsmath,amssymb,amsthm,mathtools}
\usepackage[T1]{fontenc}
\usepackage{lmodern}
\DeclareMathSizes{10}{10}{8}{8}
\DeclareMathSizes{10.95}{10.95}{8}{8}
\setlength{\emergencystretch}{1.8em}
\usepackage[hidelinks]{hyperref}
\usepackage{microtype}
\newtheorem{theorem}{Theorem}[section]
\newtheorem{lemma}[theorem]{Lemma}
\newtheorem{corollary}[theorem]{Corollary}
\newcommand{\wt}{\operatorname{wt}}
\newcommand{\coef}[2]{[z^{#1}]#2}
\newcommand{\pos}[1]{(#1)_+}
\newcommand{\Prb}{\mathbb P}
\newcommand{\NN}{\mathbb Z_{\geq0}}
\title{Rank-unimodality of one-descent composition ideals}
\hypersetup{unicode=true,
  pdftitle={Rank-unimodality of one-descent composition ideals},
  pdfauthor={Zhang Yibo and Zhang Zhanhe}}
% The byline uses the title-page space otherwise reserved for an empty date.
\makeatletter
\renewcommand{\@maketitle}{%
  \begingroup\centering
  \vspace*{0.6em}%
  {\LARGE\@title\par}\vspace{0.85em}%
  \@author\par
  \endgroup
  \vspace{0.85em}}
\makeatother
\author{\begin{tabular}{@{}c@{\hspace{2em}}c@{}}
  \normalsize Zhang Yibo &
  \normalsize Zhang Zhanhe\\[-1pt]
  \footnotesize\href{https://orcid.org/0009-0002-2005-5781}{ORCID 0009-0002-2005-5781} &
  \footnotesize\href{https://orcid.org/0009-0007-2883-9640}{ORCID 0009-0007-2883-9640}
\end{tabular}\endgraf\vspace{3pt}%
{\footnotesize Both authors contributed equally as joint first authors. Zhang Zhanhe is the corresponding author.\endgraf}}
\date{}
\begin{document}
\maketitle
\begin{abstract}
The principal ideal of a composition consists of its distinct generalized
subwords, ranked by the sum of their parts. We prove rank-unimodality for
every composition with at most one descent and for every bitonic composition.
Both the part sizes and the number of parts are unrestricted. The main
ingredient is a coupling for a monotone host. If a subword is chosen uniformly
among those of weight at most \(n\), its length and its unused weight can be
made to increase weakly when \(n\) increases by one. A decreasing threshold
on the next available part then compares the positive difference of two
ordinary rank coefficients with an earlier cumulative ratio. This supplies
the coefficient inequality in the hereditary reduction at a strict peak.
\end{abstract}

\section{Principal ideals and a coefficient inequality}
A composition is a finite word on the positive integers. Its weight is the
sum of its letters. The empty word has weight zero. Write \(u\preceq w\)
if the letters of \(u\) can be assigned in order to distinct positions of
\(w\), each assigned letter at most the letter in its position. Set
\[
                  F_w(z)=\sum_{u\preceq w}z^{\wt(u)}.
\]
Each word is counted once, independently of its number of embeddings. A
descent is an index \(i\) with \(w_i>w_{i+1}\). A bitonic word weakly
increases and then weakly decreases. A sequence is unimodal if it weakly
increases up to some index and weakly decreases from that index onwards.

\begin{theorem}\label{thm:main}
If \(w\) has at most one descent, the coefficients of \(F_w\) are unimodal.
The same holds if \(w\) is bitonic.
\end{theorem}

Under the usual correspondence between compositions and layered permutations,
weight is permutation length and generalized subword containment is pattern
containment. The theorem settles the one-descent case of Vatter's
Conjecture~2.6 \cite{vatter}; the descent here is in the sequence of layer
sizes. Sagan's rectangular theorem and the monotone-host corollary give
earlier cases \cite{sagan,vatter}. Zhang derived the symmetric peak
identity of Section 4 and proved the cases of a single final part and
of a peak at least the combined arm weight minus two \cite{zhang}.
Chen reports rank-unimodality for hosts of length at most six, tail
log-concavity for bitonic hosts with parts at most five, and tail
log-concavity for the ordinary product of two monotone arms \cite{chen}.

We call a weakly increasing host an arm. We use zero for the maximum
of an empty arm and \(\max(X,Y)\) for the largest part of either arm.
For an arm \(X\), let
\(A_i=[z^i]F_X\), \(G_i=\sum_{k=0}^i A_k\), and
\(D_i=\max(A_i-A_{i-1},0)\) for \(i\ge2\). Coefficients at negative indices
are zero. The inequality that will be used at a peak is
\begin{equation}
            G_jD_{N-1}\ \ge\ G_{j-1}D_N
            \qquad (1\le j\le N-3,\ D_{N-1}D_N>0).              \label{eq:bridge}
\end{equation}
We will obtain it also for \(j=N-2\) when \(N\ge3\).
The distinctions between \(G\), the coefficients \(A_i\), and the clipped
differences \(D_i\) are essential. Already \(F_{(1,2)}=1+z+2z^2+z^3\) fails coefficient
log-concavity at rank one.

\section{Length and unused weight for a monotone host}
Given \(b_0\ge b_1\ge\cdots\ge b_{H-1}\ge1\), set \(b_H=0\),
\(S_t=z+\cdots+z^t\), \(S_0=0\), and
\[
 P_0=1,\qquad P_\ell=\prod_{r=0}^{\ell-1}S_{b_r}\quad(1\le\ell\le H).
\]
A word of length \(\ell\) fits the decreasing host \((b_0,\ldots,b_{H-1})\)
if and only if it fits its first \(\ell\) positions. Hence
\(F=\sum_{\ell=0}^H P_\ell\) is the rank polynomial. Reversal gives the
same weight and length counts for the reversed increasing host. Write
\(p_\ell(k)=[z^k]P_\ell\),
\(C_\ell(n)=\sum_{k\le n}p_\ell(k)\), and
\(C(n)=\sum_\ell C_\ell(n)\), with \(C_\ell(n)=0\) for \(n<0\).
In particular \(C(n)>0\) for \(n\ge0\), including beyond the total weight.

\begin{theorem}\label{thm:coupling}
At budget \(n\), choose uniformly from the words of weight at most \(n\).
Write \(K\) for the weight, \(J\) for the number of parts, and \(n-K\)
for the unused weight. The choices at budgets \(n\) and \(n+1\) admit
a coupling in which all three quantities \(K,J,n-K\) increase weakly.
Moreover \(C(n)\) is log-concave for every \(n\ge0\). In particular,
for every nonincreasing sequence of nonnegative integer thresholds
\(s_0,\ldots,s_H\),
\begin{equation}
          \frac{\sum_{\ell=0}^{H}C_\ell(n-s_\ell)}{C(n)}
          \quad\hbox{is nondecreasing in }n.                       \label{eq:threshold}
\end{equation}
\end{theorem}

\begin{proof}
First condition on exact weight. For \(i<j,\ k<k'\), the layer counts obey
\begin{equation}
              p_i(k)p_j(k')\ge p_i(k')p_j(k).                       \label{eq:tp}
\end{equation}
Indeed \(P_j=P_i\,\prod_{i\le r<j}S_{b_r}\).
The coefficients of \(P_i\) form a log-concave sequence on an interval.
Each nonnegative translate has increasing ratio to that sequence on
common support. In cross-multiplied form this remains true at either
zero boundary. Sum over the translates to get \eqref{eq:tp}.
For completeness, the indicator of an integer interval has nonnegative order-two minors
in its Toeplitz matrix. This property is equivalent, on interval support,
to log-concavity, by multiplying successive ratio inequalities. For two
such matrices it passes to their product by the order-two Cauchy--Binet
formula. For a prescribed pair of rows and columns all nonzero summands
are finite. This proves the asserted log-concavity for products of the
shifted intervals, including the singleton \(P_0\).

Summing \eqref{eq:tp} over \(i<d\le j\) shows that
\(\Prb(J\ge d\mid K=k)\) is nondecreasing in \(k\), wherever that
conditional distribution is defined. All possible weights form the
full interval from zero to \(\sum b_r\). The distribution at a higher
weight can therefore be coupled with the one at a lower weight to have
at least as many parts. Explicitly, use the same uniform random number
in the two inverse distribution functions.

Next suppose the successive ratios \(C(i)/C(i-1)\), \(1\le i\le n\),
are nonincreasing. For the moment \(n\ge1\). We construct the coupling
from budget \(n-1\) to budget \(n\). At budget \(v\) the weight has
distribution function \(C(\min(k,v))/C(v)\) for \(k\ge0\).
Increasing the budget increases the weight in stochastic order. It
increases it by at most one in the common-quantile coupling, since
\[
        \frac{C(k)}{C(n-1)}\le\frac{C(k+1)}{C(n)}
                \qquad (0\le k\le n-1).
\]
For \(k=n-1\) this is equality, and for smaller \(k\) it is the
assumed ordering of ratios. Thus, in this coupling,
\(K_{n-1}\le K_n\le K_{n-1}+1\). Both the weight and the unused
weight increase weakly. Conditional on these two weights, use the
length coupling just proved. If the actual words are desired, choose
uniformly within each layer of fixed weight and length. Each marginal
then has precisely the required law. Since \(s_J\) weakly decreases as
\(J\) increases, the probability of \(n-K_n\ge s_{J_n}\) increases at
this step. That probability is exactly \eqref{eq:threshold}.

It remains to justify the assumed ordering of ratios. The layer identity
\((1-z)S_t=z-z^{t+1}\) gives, on summing \(P_{\ell+1}=P_\ell S_{b_\ell}\),
\[
 (1-2z)F=1-z-z\sum_{\ell=0}^H z^{b_\ell}P_\ell .
\]
The terminal \(b_H=0\) accounts for the last layer. Dividing by \(1-z\)
and taking coefficients at \(n+1\), for all \(n\ge0\), gives
\begin{equation}
       \frac{C(n+1)}{C(n)}
            =2-\frac{\sum_\ell C_\ell(n-b_\ell)}{C(n)}.              \label{eq:rec}
\end{equation}
There is only one ratio to consider through \(n=1\). If the ratios
are ordered through \(n\), the preceding coupling makes the fraction
on the right of \eqref{eq:rec} at budget \(n\) at least its value at
\(n-1\), since \(b_0\ge\cdots\ge b_H\). Thus the next ratio is at
most the last. This proves the induction. It includes the empty host,
the budget-zero comparison, and the constant tail of \(C\). The
coupling and \eqref{eq:threshold} now hold at every step.
\end{proof}

For equal capacities, log-concavity of \(C\) also follows by applying
the compound-distribution theorem of Bad\'ia, Sang\"uesa and Federgruen
\cite{badia} to uniform positive summands and a truncated geometric
distribution on the number of summands. In Theorem~\ref{thm:coupling},
the uniform summands can have different supports. The exact number
of words in each layer enters both the conditional law and
\eqref{eq:rec}.

\section{The derivative and the cumulative sequence of the same arm}
Put \([t]=1+\cdots+z^{t-1}\). We may write a nonrectangular arm
as \(X=1^\alpha(t_1,\ldots,t_h)\), with \(\alpha\ge1\),
\(2\le t_1\le\cdots\le t_h\), \(h\ge1\). If the minimum letter of a
nonempty arm is \(c>1\), replace its first letter by \(c\) copies of one:
the initial polynomial in both cases is \([c+1]\), and every subsequent
append applies \(F\mapsto1+S_tF\). A word of ones or a single letter
has an interval polynomial.

Form the reference host with decreasing capacities \(t_h,\ldots,t_2\),
omitting \(t_1\). Use the preceding notation \(P_\ell,C_\ell,C\)
for \(0\le\ell<h\), and set \(P_h=\prod_{i=1}^h S_{t_i}\). The two
sequences measured against \(C\) below run in opposite directions.

\begin{theorem}\label{thm:sandwich}
For every such arm the sequence \(D_{m+2}/C(m)\) is nonincreasing,
and \(G_m/C(m)\) is nondecreasing, for \(m\ge0\). The reference
\(C\) is positive and log-concave.
\end{theorem}

\begin{proof}
Start the append recurrence with \((1-z)^{-1}\) in place of
\([\alpha+1]\), and let \(a_m\) be the resulting signed coefficient
difference at ranks \(m+2,m+1\). The initial differences vanish.
An append with capacity \(t\) changes their generating function by
\(d\mapsto[t-1]+S_t d\). Iterating gives
\[
 a(z)=\sum_{\ell=0}^{h-1}[t_{h-\ell}-1]P_\ell(z),\qquad
 a_m=C(m)-\sum_{\ell=0}^{h-1}C_\ell(m-t_{h-\ell}+1).
\]
The thresholds \(t_{h-\ell}-1\) decrease with \(\ell\). By
\eqref{eq:threshold} the ratio \(a_m/C(m)\) is nonincreasing.

We give the finite-padding argument to keep track of the positive part.
In the signed difference, the correction from replacing the initial
series by \([\alpha+1]\) is the subtraction of the polynomial
\(T=z^{\alpha-1}P_h\). Hence \(D_{m+2}=(a_m-T_m)_+\).
For \(i=h-\ell\), its corresponding summand in \(a\) and the whole
correction have the same last factor:
\[
  A_i=P_\ell[t_i-1],\qquad
  T=P_\ell\, z^{\alpha+i-1}\prod_{j=1}^i[t_j].
\]
The coefficients of the second factor of \(T\) are nondecreasing
on the support \(0,\ldots,t_i-2\) of \([t_i-1]\), since the factor
\([t_i]\) has no outgoing term on that interval. Its ratio to
\([t_i-1]\) is therefore increasing in the cross-product sense, also
at the zero boundaries. Convolution with \(P_\ell\) preserves that
order: expand each resulting determinant by Cauchy--Binet and use
the nonnegative Toeplitz minors of \(P_\ell\). Thus
\(T_m(A_i)_{m'}\le T_{m'}(A_i)_m\) for \(m<m'\). Sum over \(i\).
It follows that \((a_m-T_m)_+/a_m\) is nonincreasing on the
support of \(a\). Every summand of \(a\) has interval support
starting at \(\ell\) and reaching at least \(\ell\); the summands
together have interval support from zero. Beyond it \(D_{m+2}=0\).
Multiplying the two decreasing nonnegative ratios proves the first
assertion with its zero tail.

For the second assertion the actual cumulative is exactly
\[
        G_m=C(m)+[z^m]\,\frac{P_h[\alpha+1]}{1-z}.
\]
Write the additional coefficient as \(V_m\). For each
\(\ell<h\), \(V\) is a nonnegative convolution of \(C_\ell\):
the remaining factor is \([\alpha+1]\prod_{i=1}^{h-\ell}S_{t_i}\).
The sequence \(C_\ell\) is log-concave as a convolution of
\(P_\ell\) with the all-ones sequence. Its translates therefore
give \(V_{m+1}C_\ell(m)\ge V_m C_\ell(m+1)\), with zero extension
at negative indices. Summing gives the same inequality with \(C\)
in place of \(C_\ell\), and hence \(G/C\) is nondecreasing.
Log-concavity of \(C\) is Theorem~\ref{thm:coupling}.
\end{proof}

\begin{corollary}\label{cor:full}
For every normalized nonrectangular arm, every \(N\ge3\), and
every \(1\le j\le N-2\), \(G_jD_{N-1}\ge G_{j-1}D_N\).
\end{corollary}
\begin{proof}
For a positive adjacent derivative pair, apply Theorem~\ref{thm:sandwich}
and compare the adjacent ratios of the positive log-concave sequence \(C\):
\[
 \frac{D_N}{D_{N-1}}\le \frac{C(N-2)}{C(N-3)}
       \le\frac{C(j)}{C(j-1)}\le\frac{G_j}{G_{j-1}}.
\]
A zero current derivative gives the inequality immediately. A zero
previous derivative forces a zero current derivative by the first
assertion of Theorem~\ref{thm:sandwich}.
\end{proof}


\section{Passage to a peak}
We record the reduction from single-arm inequalities to the coefficients
at a peak. Its two inductions concern different objects. The first concerns
the ordinary product for shorter arms; the second concerns rank polynomials
at smaller peaks and at proper deletions of the same peak. The technical
companion to this article contains the coefficient and boundary lemmas
for these inductions \cite{companion}.

\begin{lemma}[Hereditary coefficient criterion]\label{lem:criterion}
Let \(\mathcal A\) be a class of increasing normalized arms closed under
deleting the final part, with its interval and empty bases included. Suppose
every nonrectangular arm in \(\mathcal A\) satisfies \eqref{eq:bridge}.
For \(X,Y\in\mathcal A\), write \(A=F_X\), \(B=F_Y\),
\(q_T=z^{\wt(T)}F_T(1/z)\), and \(p=q_Xq_Y\). Then
\(\bigl(2^j(p_j-p_{j-1})_+\bigr)_{j\ge0}\) is unimodal.
For every integer \(M>\max(X,Y)\) the three polynomials
\begin{equation}
      A+B+(S_M-1)AB,\qquad A+(S_M-1)AB,\qquad B+(S_M-1)AB
                                                                  \label{eq:peak}
\end{equation}
have unimodal coefficients.
\end{lemma}

We describe the use of the hypothesis; this identifies the precise
input supplied by Corollary~\ref{cor:full}. Set \(P=AB\),
\(e_n=2P_{n-1}-P_n\), \(E_B=1-z+(2z-1)B\), and \(w_j=[z^j](1-z)E_B\).
While both ordinary factors rise, the single-arm coefficient bounds
give \(w_j\ge0\) on the prefix in question, with \(w_j=0\) for \(j<3\).
Write \(J_n=[z^n]E_BA=\sum_{j=3}^nw_jG_{n-j}\), where \(G\) is the
cumulative sequence of \(A\). Termwise application of \eqref{eq:bridge}
gives
\[
        J_{n+1}D_n\ge J_nD_{n+1},\qquad J_{n+1}\ge J_n,
        \qquad e_n=J_n-D_n .
\]
The unused term in the first sum is nonnegative. If \(e_{n+1}\ge0\)
and both differences are positive, then
\[
 J_{n+1}(e_{n+1}-e_n)
    =J_{n+1}D_n-J_nD_{n+1}
      +(J_{n+1}-J_n)e_{n+1}\ge0.
\]
Here \(J_{n+1}>0\); if \(e_{n+1}<0\), the ratio inequality instead gives
\(J_n<D_n\). This is the comparison on the common rising prefix, including
the case \(e_{n+1}=0\). The companion proves the mixed and nonrising
comparisons for true arms without hypothesis \eqref{eq:bridge}.

For clarity, the rest of the criterion uses the following identities.
Deleting \(t\) from a heavier arm \(X=Zt\) yields
\[
           q_Xq_Y=[t]q_Zq_Y+z^{\wt(X)}q_Y .
\]
The closed interval comparison of the companion propagates the weighted
positive-difference conclusion from the shorter product to its interval
convolution. The terminal summand begins after the first ordinary mode
of \(q_Xq_Y\), so the new positive differences are a truncation of the
window differences. This proves the first conclusion by induction with
the rectangular and one-part bases. In the peak induction the reversal
of \(1+(S_M-1)B\) is
\[
               \delta_M(Y)=([M]-z^M)q_Y+z^{\wt(Y)+M}.
\]
Reversal of the corresponding relative in \eqref{eq:peak} is
\(q_X\delta_M(Y)\). At a least failing peak, choose a relative with
the fewest parts. Deletion from the heavier arm gives the already
unimodal proper relative; its interval convolution supplies the required
initial signs. The weighted positive-difference conclusion controls
the remaining central signs. Stripping all copies of the largest arm
part gives the lower-peak initial case. The reflection and terminal
joins give both relatives and the complete core. The ranges at the
first interval boundary, at a string of zero differences, and at the
transition to the lower peak are included in the companion's reduction.
In particular this criterion uses \eqref{eq:bridge} at all deletion
prefixes. No assertion about an individual tested peak enters it.

\begin{proof}[Proof of Theorem~\ref{thm:main}]
Take \(\mathcal A\) to be all increasing arms, normalized as in
Section~3. Each final deletion is again an increasing arm. All
nonrectangular members, including every proper prefix, satisfy
\eqref{eq:bridge} by Corollary~\ref{cor:full}. Lemma~\ref{lem:criterion}
therefore gives all three conclusions in \eqref{eq:peak} at every
strict peak.

To identify these polynomials, take any words \(U,V\) with letters at
most \(M\), and write \(A=F_U,\ B=F_V\). A word under \(U,M,V\) either
fits \(U\), or has a unique longest prefix \(u\) fitting \(U\) followed
by a letter \(x\le M\) for which \(ux\) fails in \(U\). The remainder
fits \(V\): in any embedding the failing letter lies at or to the right
of the displayed \(M\). Conversely \(u,x\), and any word fitting \(V\)
give a valid word by assigning \(x\) to \(M\). The possible failing
extensions have polynomial \(S_M A-(A-1)\). Thus
\[
            F_{U,M,V}=A+\{1+(S_M-1)A\}B.
\]
This proves the first identity in \eqref{eq:peak} for the core and
also shows that reversing either arm leaves its polynomial unchanged.

For a word with exactly one descent, let \(M\) be the upper letter of that
descent. Write the word \(X,M^{r+1},Y,E\), where \(X,Y,E\) are weakly
increasing, all parts of \(X,Y\) are strictly below \(M\), and every
part of \(E\) is at least \(M\). Start with the strict core \(X,M,Y\).
Each extra central \(M\) sends its polynomial \(F\) to \(1+S_MF\),
as follows from the preceding identity and \(F_{X,M}=1+S_MF_X\).
Each letter \(t\) of \(E\) is a current global maximum and likewise
sends \(F\) to \(1+S_tF\). If \(F\) has constant term one and is
unimodal, interval convolution preserves unimodality, and the first
two coefficients of \(1+S_tF\) are \(1,1\). Thus each step preserves
the conclusion. A monotone host follows directly by the same append
induction. For a bitonic host choose its central maximum, strip
repeated maxima, and reverse its decreasing right arm. Reversal
preserves \(F_Y\), so the strict-core result and the same restoration
complete the proof.
\end{proof}

\section{Scope and verification}
The distinction between a cumulative statement and a rank statement can
be seen even for \(1,3,2,3\). Its cumulative sequence \(1,4,6,9\) is
log-concave, while the coefficients have a strict decrease followed by a
strict increase. In the proof of Theorem~\ref{thm:main},
Theorem~\ref{thm:coupling} supplies a joint monotonicity statement;
the specific thresholds \(t_{h-\ell}-1\) identify an ordinary difference.
The separate deepest-layer and finite-padding comparisons are what turn
that statement into \eqref{eq:bridge}.

The companion includes the Lean proofs of the finite transport,
the induction in \eqref{eq:rec}, the exact arm recurrences (including
finite padding), and the scalar ratio transfers. It also
includes exact checks that construct distinct subwords by their first
possible embedding, independently of the peak formula. The identification
of the composition layers, the application of the transport to them,
and the full strict-peak reduction are mathematical proofs. General
hosts with multiple descents outside the bitonic class remain covered
only to the extent of other known results.

\begin{thebibliography}{9}
\bibitem{badia} F.~G. Bad\'ia, C. Sang\"uesa and A. Federgruen,
Log-concavity of compound distributions with applications in operational
and actuarial models, \emph{Probability in the Engineering and
Informational Sciences} \textbf{35} (2021), no.~2.
\url{https://doi.org/10.1017/S0269964819000329}.
\bibitem{chen} L. Chen, Rank-unimodality and tail log-concavity of
composition ideals, preprint, 2026. \url{https://ssrn.com/abstract=7385218}.
\bibitem{sagan} B.~E. Sagan, Compositions inside a rectangle and
unimodality, \emph{Journal of Algebraic Combinatorics} \textbf{29}
(2009), 405--411. \url{https://arxiv.org/abs/0707.1052}.
\bibitem{vatter} V. Vatter, An assortment of problems in permutation
patterns: unimodality, equivalence, derangements, and sorting,
arXiv:2602.16355v3 (2026). \url{https://arxiv.org/abs/2602.16355}.
\bibitem{zhang} Z. Zhang, Rank-unimodality for one-descent families in
generalized subword order, preprint, 2026. \url{https://ssrn.com/abstract=7397439}.
\bibitem{companion} Technical companion to the present paper,
\emph{The hereditary strict-peak reduction}, with mathematical proofs,
finite certificates, and verification sources. Included with the manuscript.
\end{thebibliography}
\end{document}